% Gerado automaticamente.
\begin{table}[!ht]
\caption{Species whose strength-based class assignment disagrees with the class implied by their measured stiffness.\label{tab:discordancia}}
\begin{threeparttable}
\begin{tabular*}{\columnwidth}{@{\extracolsep{\fill}}llrrrr@{\extracolsep{\fill}}}
\toprule
Species & Class by & $E_{M0}$ & Class by & Shift & Deviation \\
 & $f_{c0,k}$ & (MPa) & $E_{M0}$ & (classes) & vs. class \\
\midrule
Angelim-pedra & D40 & 10755 & D20 & -2 & -25,8\% \\
Piolho & D40 & 11790 & D20 & -2 & -18,7\% \\
Quarubarana & D20 & 8842 & <D20 & -1 & -11,6\% \\
Cedro-doce & D20 & 8866 & <D20 & -1 & -11,3\% \\
Umirana & D30 & 10794 & D20 & -1 & -10,1\% \\
Cedro-amargo & D20 & 9176 & <D20 & -1 & -8,2\% \\
Angelim-pedra-verdadeiro & D50 & 15215 & D40 & -1 & -7,8\% \\
Castelo & D30 & 11375 & D20 & -1 & -5,2\% \\
Angelim-araroba & D30 & 11457 & D20 & -1 & -4,5\% \\
Guaicara & D40 & 13866 & D30 & -1 & -4,4\% \\
Cafearana & D40 & 13982 & D30 & -1 & -3,6\% \\
Rabo-de-arraia & D40 & 14324 & D30 & -1 & -1,2\% \\
Oiticica-amarela & D40 & 14491 & D30 & -1 & -0,1\% \\
Angico-preto & D50 & 16498 & D40 & -1 & -0,0\% \\
Itauba & D40 & 17345 & D50 & +1 & 19,6\% \\
Angelim-ferro & D50 & 19938 & D60 & +1 & 20,8\% \\
Angelim-saia & D40 & 17561 & D50 & +1 & 21,1\% \\
Catanudo & D30 & 14729 & D40 & +1 & 22,7\% \\
Canafistula & D30 & 14769 & D40 & +1 & 23,1\% \\
Mandioqueira & D40 & 18679 & D50 & +1 & 28,8\% \\
Branquilho & D30 & 15490 & D40 & +1 & 29,1\% \\
Casca-grossa & D30 & 16802 & D50 & +2 & 40,0\% \\
Goiabao & D30 & 18367 & D50 & +2 & 53,1\% \\
\bottomrule
\end{tabular*}
\begin{tablenotes}\footnotesize
\item 23 of the 40 species (58\%) fall in a different class when assigned by measured stiffness rather than by characteristic compressive strength. 14 species (35\%) are stiffer on paper than in reality, so class-based design is unconservative for the deflection limit state.
\end{tablenotes}
\end{threeparttable}
\end{table}
